\chapter{Mengingat: dari RNN ke LSTM dan Transformer}
\label{ch:lstm}

Di semester pertama, kuliah LSTM diberikan langsung oleh Sepp Hochreiter, yang menganalisis
masalah vanishing gradient pada 1991 dan bersama Jürgen Schmidhuber menerbitkan LSTM pada 1997.
Kuliahnya tidak berhenti di sejarah. Bagian akhirnya membahas xLSTM, versi modern yang
dikembangkan di Linz untuk bersaing dengan transformer. Setahun kemudian, kuliah NLP melengkapi
sisi lain ceritanya: embedding kata dan model bahasa besar.

Banyak data datang berurutan: kata dalam kalimat, detak jantung, harga saham, langkah waktu
sebuah simulasi. Untuk data seperti ini, model harus punya ingatan. Bab ini mengikuti evolusi
ingatan buatan, dari jaringan rekuren sederhana yang cepat lupa, ke LSTM yang menjaga ingatan
dengan gerbang, ke transformer yang melihat seluruh masa lalu sekaligus, lalu kembali lagi ke
ingatan rekuren versi modern.

\section{Jaringan rekuren}

Jaringan feed-forward memproses setiap input sendiri-sendiri. \istilah{Recurrent neural network}
(RNN) membawa sebuah keadaan tersembunyi $\vh_t$ dari satu langkah ke langkah berikutnya:
\begin{equation}
\vh_t=f\big(\mW\vx_t+\bm R\,\vh_{t-1}+\bm b\big),\qquad \hat{\vy}_t=g(\bm V\vh_t).
\end{equation}
Matriks $\bm R$ adalah bobot rekuren. Bobot yang sama dipakai di setiap langkah waktu, persis
seperti kernel konvolusi dipakai di setiap posisi gambar. Saat membaca novel, kita tidak
mengingat setiap kata. Kita membawa ringkasan berjalan di kepala: siapa tokohnya, apa masalahnya,
apa yang baru saja terjadi. Setiap kalimat baru memperbarui ringkasan itu. Keadaan tersembunyi
$\vh_t$ adalah ringkasan berjalan tersebut.

Kuliah membahas beberapa varian klasik. Jaringan Jordan mengumpankan kembali keluarannya,
jaringan Elman mengumpankan kembali lapisan tersembunyinya, jaringan rekuren penuh menghubungkan
semua unit, jaringan NARX memakai input dan keluaran yang tertunda, dan \istilah{time-delay neural
network} hanya melihat jendela waktu yang tetap.

Untuk melatih RNN, \istilah{backpropagation through time} (BPTT) membuka jaringan menjadi jaringan
feed-forward yang dalam, satu lapisan untuk setiap langkah waktu, lalu menjalankan backprop
biasa. Karena bobotnya dipakai bersama, gradien untuk $\bm R$ adalah jumlah kontribusi dari semua
langkah. Untuk deret yang panjang, BPTT dipotong setelah $k$ langkah supaya memori dan waktu
tetap terbatas. Alternatifnya adalah \istilah{real-time recurrent learning} (RTRL), yang membawa
turunan keadaan terhadap semua bobot maju bersama waktu. Gradien tersedia di setiap langkah tanpa
menyimpan masa lalu, tetapi harganya mahal: untuk jaringan rekuren penuh dengan $n$ unit,
biayanya tumbuh seperti $n^4$ per langkah. Dalam bahasa \cref{ch:dlbasic}, RTRL adalah forward
mode, sedangkan BPTT adalah reverse mode.

\section{Gradien yang hilang dan meledak}

Masalah berikut mengubah sejarah jaringan saraf. Delta di langkah $t-k$ bergantung pada delta di
langkah $t$ melalui perkalian $k$ buah Jacobian:
\begin{equation}
\frac{\partial\vh_t}{\partial\vh_{t-k}}=\prod_{\tau=t-k+1}^{t}\frac{\partial\vh_\tau}{\partial\vh_{\tau-1}}
=\prod_{\tau}\mathrm{diag}\big(f'(\bm s_\tau)\big)\,\bm R.
\end{equation}
Kalau norma setiap faktor lebih kecil dari satu, hasil kalinya menyusut secara eksponensial.
Inilah \istilah{vanishing gradient}. Kalau lebih besar dari satu, ia meledak. Hochreiter
menganalisis gejala ini pada 1991 dalam tesis diplomanya \citep{hochreiter1991}, dan
\citet{bengio1994} menunjukkan secara independen bahwa ketergantungan jangka panjang memang sulit
dipelajari dengan gradient descent.

Seberapa cepat gradien hilang? Ambil RNN skalar $h_t=\tanh(r\,h_{t-1}+w\,x_t)$ dengan $r=0{,}5$.
Karena turunan tanh tidak pernah melebihi satu, setiap langkah mengalikan gradien dengan paling
banyak $0{,}5$. Setelah dua puluh langkah, faktornya paling besar $0{,}5^{20}$, kurang dari satu
per sejuta. Informasi dari dua puluh langkah sebelumnya praktis tidak lagi memengaruhi
perubahan bobot. Dengan $r=1{,}5$ dan unit yang bekerja di daerah linearnya, faktor yang sama bisa
mencapai sekitar 3300.

Gradien yang meledak relatif mudah ditangani dengan \istilah{gradient clipping}: kalau norma gradien
melewati batas tertentu, skalakan turun. Gradien yang hilang lebih sulit, karena tidak ada sinyal
yang bisa diperbesar. Solusinya harus datang dari arsitektur.

\section{LSTM}

\citet{hochreiter1997} memecahkan masalah ini dengan dua gagasan yang di kuliah disebut sebagai inti
LSTM: \istilah{constant error carousel} dan \istilah{gating}.

Gagasan pertama sederhana. Bayangkan satu unit ingatan $c_t$ yang diperbarui dengan penjumlahan,
bukan perkalian, yaitu $c_t=c_{t-1}+(\text{sesuatu yang baru})$. Turunannya
$\partial c_t/\partial c_{t-1}$ bernilai tepat satu. Error yang mengalir mundur lewat $c$ tidak
menyusut dan tidak membesar. Ia berputar tanpa berubah, seperti penumpang di komidi putar.

Penjumlahan tanpa kendali akan mengisi ingatan dengan sampah. Karena itu LSTM memberi gerbang
yang dipelajari. Gerbang input memutuskan apa yang boleh masuk, dan gerbang output memutuskan apa
yang boleh dibaca. Gerbang lupa (\istilah{forget gate}) ditambahkan kemudian oleh \citet{gers2000}
agar sel bisa mengosongkan dirinya. Versi yang kini dipakai di mana-mana berbentuk
\begin{equation}
\begin{aligned}
\vz_t&=\tanh(\mW_z\vx_t+\bm R_z\vh_{t-1}+\bm b_z), &
\bm i_t&=\sigma(\mW_i\vx_t+\bm R_i\vh_{t-1}+\bm b_i),\\
\bm f_t&=\sigma(\mW_f\vx_t+\bm R_f\vh_{t-1}+\bm b_f), &
\bm o_t&=\sigma(\mW_o\vx_t+\bm R_o\vh_{t-1}+\bm b_o),\\
\bm c_t&=\bm f_t\odot\bm c_{t-1}+\bm i_t\odot\vz_t, &
\vh_t&=\bm o_t\odot\tanh(\bm c_t).
\end{aligned}
\end{equation}
Di sini $\vz_t$ adalah input sel, $\bm i_t$, $\bm f_t$, dan $\bm o_t$ adalah gerbang input, lupa,
dan output, dan $\bm c_t$ adalah keadaan sel. Sepanjang jalur sel, $\partial\bm c_t/\partial\bm c_{t-1}=\bm f_t$.
Selama gerbang lupa dekat satu, error mengalir hampir tanpa hambatan, dan jaringan sendiri yang
memutuskan kapan harus melupakan.

Saya membayangkan keadaan sel sebagai buku catatan. Gerbang input adalah stabilo: hanya hal yang
ditandai yang ditulis ke buku. Gerbang lupa adalah penghapus, karena kadang halaman lama perlu
dibersihkan. Gerbang output adalah keputusan tentang halaman mana yang dibacakan sekarang.
Tulisan di buku tidak memudar dengan sendirinya. Ia hanya berubah kalau ditulis atau dihapus.

Gambaran ini juga membantu menghindari beberapa salah paham. LSTM tidak menghilangkan semua
masalah gradien. Jalur sel memang aman, tetapi jalur yang melewati gerbang dan $\vh$ masih bisa
mengalami vanishing. LSTM membuat belajar jangka panjang menjadi mungkin, bukan otomatis. Bias
gerbang lupa juga sering diinisialisasi positif, misalnya satu, supaya di awal pelatihan sel
cenderung mengingat. Dan pembagian peran tanh dan sigmoid bukan kebetulan: gerbang memakai sigmoid
karena nilainya harus berada di antara nol dan satu, seperti keran, sedangkan input sel memakai
tanh karena isinya boleh positif atau negatif.

LSTM mendominasi pemrosesan deret selama dua dekade, dari pengenalan suara di ponsel sampai
terjemahan mesin. Kuliah juga menampilkan contoh di luar bahasa, seperti klasifikasi protein dari
urutan asam amino dan analisis interval detak jantung.

\section{Attention dan transformer}

LSTM memampatkan seluruh masa lalu ke dalam vektor berukuran tetap. Untuk menerjemahkan kalimat
yang panjang, pemampatan ini menjadi leher botol. \citet{bahdanau2015} membiarkan dekoder melihat
kembali semua keadaan enkoder dan memilih yang relevan. Mekanisme ini disebut
\istilah{attention}. Transformer \citep{vaswani2017} kemudian membuang rekurensi sama sekali dan hanya
memakai attention:
\begin{equation}
\mathrm{Attention}(\bm Q,\bm K,\bm V)=\mathrm{softmax}\!\Big(\frac{\bm Q\bm K^\top}{\sqrt{d_k}}\Big)\bm V,
\qquad \bm Q=\mX\mW_Q,\ \bm K=\mX\mW_K,\ \bm V=\mX\mW_V.
\end{equation}
Baris ke-$i$ hasilnya adalah rata-rata berbobot dari semua \istilah{value}, dengan bobot dari
kemiripan \istilah{query} ke-$i$ dengan setiap \istilah{key}.

Perpustakaan adalah gambaran yang pas. Anda datang dengan sebuah pertanyaan, yaitu query. Setiap
buku punya label di punggungnya, yaitu key, dan isi, yaitu value. Anda membandingkan pertanyaan
dengan semua label, lalu membaca isi buku dengan perhatian yang sebanding dengan kecocokannya.
Buku yang labelnya cocok sekali dibaca serius, yang tidak cocok hanya dilirik. Jawaban Anda
adalah campuran dari isi buku-buku itu.

Sebagai contoh kecil, ambil query $\bm q=(1,0)$, dua key $\bm k_1=(1,0)$ dan $\bm k_2=(0,1)$, value
$v_1=10$ dan $v_2=0$, serta $d_k=2$. Skornya $\bm q\cdot\bm k_1/\sqrt2\approx0{,}707$ dan $0$.
Softmax mengubahnya menjadi bobot $0{,}670$ dan $0{,}330$, karena $e^{0{,}707}\approx2{,}028$
dibagi $2{,}028+1$. Keluarannya $6{,}70$. Pembagian dengan $\sqrt{d_k}$ menjaga skor tidak terlalu
besar saat dimensinya tinggi, supaya softmax tidak jenuh menjadi pilihan tunggal yang keras.

Mengapa tepat $\sqrt{d_k}$? Anggap komponen $q_i$ dan $k_i$ saling independen dengan rata-rata nol dan
varians satu. Hasil kali titik $\bm q\cdot\bm k=\sum_{i=1}^{d_k}q_ik_i$ adalah jumlah $d_k$ suku yang
masing-masing punya rata-rata nol dan varians $\E[q_i^2]\E[k_i^2]=1$, sehingga variansnya $d_k$ dan
simpangan bakunya $\sqrt{d_k}$. Untuk $d_k=64$, skor tipikal berukuran sekitar delapan. Selisih delapan
pada input softmax berarti rasio bobot $e^8\approx3000$, sehingga hampir semua bobot jatuh ke satu kunci
dan gradien ke kunci-kunci lain praktis nol. Membagi dengan $\sqrt{d_k}$ mengembalikan simpangan baku skor
ke satu, berapa pun dimensinya.

Satu blok transformer berisi \istilah{multi-head attention}, yaitu beberapa attention paralel dengan
proyeksi yang berbeda yang hasilnya digabung, lalu jaringan feed-forward yang bekerja per posisi.
Keduanya dibungkus koneksi residual dan layer normalization. Karena attention tidak peduli urutan,
informasi posisi ditambahkan lewat \istilah{positional encoding}. Untuk model bahasa yang menulis dari
kiri ke kanan, sebuah \istilah{causal mask} mencegah setiap posisi melihat masa depan.

Kelebihan transformer adalah semua posisi diproses paralel selama pelatihan, dan dua token
mana pun hanya berjarak satu langkah attention. Kekurangannya, yang ditekankan di kuliah, adalah
biaya attention yang tumbuh kuadratik terhadap panjang deret, serta keharusan menyimpan key dan
value dari semua token sebelumnya ketika model membangkitkan teks.

\citet{ramsauer2021}, juga dari Linz, menunjukkan bahwa aturan update jaringan Hopfield modern
dengan keadaan kontinu,
\begin{equation}
\bm\xi^{\mathrm{baru}}=\mX\,\mathrm{softmax}\big(\beta\,\mX^\top\bm\xi\big),
\end{equation}
sama dengan attention. Kolom-kolom $\mX$ adalah pola yang disimpan dan $\bm\xi$ adalah query.
Biasanya satu langkah update sudah cukup untuk mengambil pola yang paling mirip. Dari sudut
pandang ini, attention adalah cara membaca dari sebuah ingatan asosiatif yang kapasitasnya sangat
besar.

\section{Dari bahasa ke vektor}

Kuliah NLP mengisi sisi datanya. Kata harus diubah menjadi angka sebelum masuk ke jaringan. Cara
paling sederhana adalah \istilah{bag of words}, yang menghitung frekuensi kata dalam dokumen dan
mengabaikan urutannya. Model \istilah{n-gram} memprediksi kata berikutnya dari $n-1$ kata sebelumnya
dengan menghitung frekuensi di sebuah korpus, tetapi jumlah konteks yang mungkin meledak seiring
bertambahnya $n$.

\istilah{Word embedding} memetakan setiap kata ke vektor padat berdimensi beberapa ratus, sedemikian
sehingga kata-kata yang muncul dalam konteks serupa punya vektor yang berdekatan
\citep{mikolov2013}. Hubungan seperti ``raja dikurangi pria ditambah wanita kira-kira sama dengan
ratu'' muncul dengan sendirinya. Embedding juga menyerap bias dari teks pelatihannya
\citep{bolukbasi2016}, tema yang kembali di \cref{ch:manusia}.

Model bahasa besar menggabungkan semua ini. BERT \citep{devlin2019} dilatih menebak kata yang
ditutup, sehingga ia melihat konteks kiri dan kanan. GPT-3 \citep{brown2020} dilatih menebak kata
berikutnya, dan menunjukkan bahwa dengan skala yang cukup, sebuah model bisa mengerjakan tugas
baru hanya dari beberapa contoh di dalam prompt.

\section{Rekurensi yang kembali}

Biaya kuadratik transformer mendorong pencarian arsitektur dengan biaya linear. \istilah{State space
model} memakai rekurensi linear
\begin{equation}
\vh_t=\bar{\mA}\,\vh_{t-1}+\bar{\bm B}\,x_t,\qquad y_t=\bm C\,\vh_t,
\end{equation}
yang bisa dihitung paralel justru karena linear. Mamba \citep{gu2023} membuat $\bar{\bm B}$,
$\bm C$, dan langkah diskretisasinya bergantung pada input, sehingga model bisa memutuskan apa
yang disimpan dan apa yang diabaikan, mirip gerbang LSTM. Rekurensi ini tidak lain adalah
diskretisasi persamaan diferensial biasa yang linear, $\dot{\vh}=\mA\vh+\bm Bx$, sebuah objek yang
akan kita temui lagi di \cref{ch:mlat,ch:penemuan}.

Kuliah ditutup dengan xLSTM \citep{beck2024}, yang berangkat dari tiga keterbatasan LSTM yang
disebut terang-terangan di slide. Keterbatasan pertama adalah ketidakmampuan merevisi keputusan
penyimpanan. Kalau vektor yang lebih cocok datang belakangan, LSTM sulit menimpa isi lama.
Jawabannya adalah \istilah{exponential gating}. Gerbang input $\bm i_t=\exp(\tilde{\bm i}_t)$ bisa
bernilai besar sehingga isi baru mendominasi, sementara sebuah keadaan normalisasi
$\bm n_t=\bm f_t\odot\bm n_{t-1}+\bm i_t$ menjaga skalanya, dan keluaran membaca $\bm c_t/\bm n_t$.
Agar eksponensial tidak meluap, dipakai keadaan penstabil $m_t=\max(\log f_t+m_{t-1},\log i_t)$.
Varian ini disebut sLSTM.

Keterbatasan kedua adalah kapasitas penyimpanan, karena satu sel LSTM hanya menyimpan skalar.
mLSTM mengganti sel dengan sebuah matriks dan memakai aturan update kovarians:
\begin{equation}
\bm C_t=f_t\,\bm C_{t-1}+i_t\,\bm v_t\bm k_t^\top,\qquad
\bm n_t=f_t\,\bm n_{t-1}+i_t\,\bm k_t,\qquad
\tilde{\vh}_t=\frac{\bm C_t\bm q_t}{\max\{|\bm n_t^\top\bm q_t|,1\}}.
\end{equation}
Bentuknya mirip ingatan asosiatif: simpan pasangan key dan value, lalu ambil dengan query.
Keterbatasan ketiga adalah pelatihan yang tidak bisa diparalelkan. Pada mLSTM, gerbang, key,
value, dan query tidak bergantung pada $\vh_{t-1}$, sehingga pelatihannya bisa paralel seperti
transformer. Blok sLSTM dan mLSTM kemudian ditumpuk dengan koneksi residual. Hasil yang
dilaporkan menunjukkan xLSTM bersaing dengan transformer dan state space model dalam pemodelan
bahasa, dengan biaya inferensi yang linear terhadap panjang deret.

\section{Deret di luar bahasa}

Model deret yang lahir dari bahasa dipakai jauh di luar teks. Kuliah LSTM sendiri sudah menunjukkan
klasifikasi protein dari urutan asam amino dan analisis detak jantung. Dalam genomik, urutan DNA dibaca
oleh model yang sama dengan model bahasa, dan base calling pada pengurutan nanopore memakai jaringan
rekuren untuk mengubah sinyal arus listrik menjadi huruf (\cref{ch:genomik}). Dalam kimia, string SMILES
dibangkitkan huruf demi huruf oleh LSTM untuk merancang molekul baru (\cref{ch:hayati}). Dalam sensor yang
dikenakan, deret akselerometer dikenali sebagai aktivitas (\cref{ch:pervasif}), dan di sistem
rekomendasi, urutan lagu yang baru didengar menjadi konteks untuk lagu berikutnya
(\cref{ch:rekomendasi}).

Simulasi yang berjalan dalam waktu juga sebuah deret: keadaan medan di langkah $t$ menentukan
keadaan di langkah $t+1$. Model deret dipakai di sini dengan beberapa cara. Yang paling umum
adalah memproyeksikan medan aliran ke beberapa mode POD (\cref{ch:penemuan}), lalu membiarkan LSTM
atau transformer mempelajari evolusi koefisien mode-mode itu, pendekatan yang dibahas dalam ulasan
\citet{brunton2020}. Ingatan juga berguna ketika sensor hanya mengukur sebagian keadaan, karena
masa lalu membantu menebak bagian yang tersembunyi. Transformer bahkan tidak harus berjalan
sepanjang waktu. Ia bisa memperlakukan titik-titik mesh atau partikel sebagai token. Universal
Physics Transformers \citep{alkin2024}, dari kelompok riset di Linz, memakai gagasan ini untuk
neural operator yang tidak terikat pada grid (\cref{ch:operator}). Di semua pendekatan itu,
masalah yang selalu muncul adalah galat yang menumpuk ketika model menebak berulang-ulang dari
tebakannya sendiri. Masalah ini kita bahas di \cref{ch:operator}.

\sumberbab{Bab ini mengikuti slide LSTM and Recurrent Neural Nets dari Sepp Hochreiter (RNN,
BPTT, RTRL, vanishing gradient, LSTM, attention, Mamba, xLSTM) dan slide Natural Language
Processing (representasi kata, n-gram, embedding, model bahasa besar). Bacaan utama:
\citet{hochreiter1997}, \citet{vaswani2017}, \citet{beck2024}, dan \citet{gu2023}.}
